\Needspace{8\baselineskip}\originalchapter{习题课: 弱衰变、味混合与Higgs产生}
本章收录Recitation 13--15。计算链由弱作用矩阵元进入部分宽度，再由可达末态求总寿命；CKM幺正性则控制味求和中的抵消。源题的2020课程数值保留其版本，不混作实时粒子数据。
\originalsection{轻子与赝标量衰变}
\begin{exercise}\label{ex:rec-13-1}
Recitation 13题1。已知无质量末态近似下$\tau_\mu=192\pi^3\hbar/(G_F^2m_\mu^5)$，用它求tau粒子的寿命，与实验值比较并解释差异。原件给$m_\mu=105.66$MeV、$m_\tau=1777$MeV、$\tau_\mu=2.197\times10^{-6}$s。
\end{exercise}
\begin{solution}
官方解答重编。单个$\tau\to\ell\bar\nu_\ell\nu_\tau$道的宽度由$m^5$标度给
$\Gamma_{\tau\ell}/\Gamma_\mu=(m_\tau/m_\mu)^5$。电子和muon都可产生，忽略$m_\mu/m_\tau$后的两个部分宽度相等，所以只有两个轻子道的假想总寿命
\begin{equation}
 \tau_{\rm lep}=\frac12\left(\frac{m_\mu}{m_\tau}\right)^5\tau_\mu
 \simeq8.16\times10^{-13}\mathrm{s}.
\end{equation}
源实测总寿命$2.91\times10^{-13}$s较短，因为还有强子末态。源答案取$B_e=17.8\%,B_\mu=17.4\%$，$B_{\rm lep}=0.352$；由$B_{\rm lep}=\Gamma_{\rm lep}/\Gamma_{\rm tot}$得
$\tau_{\rm tot}=B_{\rm lep}\tau_{\rm lep}\simeq2.87\times10^{-13}$s，接近该课程值。
这里$8.16\times10^{-13}$s不是单个轻子道的倒宽度，不能再漏一次因子2。非零muon质量和辐射修正解释余下偏差。原题的muon三体积分见习题\ref{ex:8701-4-1}；本题增加可达末态计数和实际总寿命比较，不重复计算同一三体积分题。
\end{solution}
\begin{exercise}\label{ex:rec-13-2}
Recitation 13题2。求$\Gamma(K^-\to e^-\bar\nu_e)/\Gamma(K^-\to\mu^-\bar\nu_\mu)$，并比较观察到的分支比。
\end{exercise}
\begin{solution}
官方pion类比扩成完整比值。对赝标量母粒子，轴流矩阵元$\langle0|\bar u\gamma^\mu\gamma^5s|K(p)\rangle=\ii f_Kp^\mu$；利用轻子Dirac方程，振幅正比$m_\ell$。两体相空间再给
\begin{equation}
 \Gamma(K\to\ell\nu)=\frac{G_F^2|V_{us}|^2 f_K^2m_Km_\ell^2}{8\pi}
 \left(1-\frac{m_\ell^2}{m_K^2}\right)^2.
\end{equation}
共同的$f_K,V_{us}$等消掉，
$R_K^{(0)}=(m_e/m_\mu)^2[(m_K^2-m_e^2)/(m_K^2-m_\mu^2)]^2\simeq2.57\times10^{-5}$。
轻电子的相空间虽大，却有更强的螺旋度抑制。原课程引用的观察值约$2.44\times10^{-5}$；须注明其源数据版本及是否包含辐射光子。电磁辐射修正会改变树级值，不能把几百分点差异当作弱耦合不普适的证据，亦不把源数值冒称最新汇总。
\end{solution}
\originalsection{幺正性与GIM抵消}
\begin{exercise}\label{ex:rec-14-1}
Recitation 14题1。证明标准CKM参数化对任意实$\theta_{12},\theta_{23},\theta_{13},\delta$均为幺正矩阵，例如$(VV^\dagger)_{11}=1$、$(VV^\dagger)_{12}=0$。记$c_{ij}=\cos\theta_{ij},s_{ij}=\sin\theta_{ij}$，原矩阵为
\begin{equation}
 V=\begin{pmatrix}
 c_{12}c_{13}&s_{12}c_{13}&s_{13}e^{-\ii\delta}\\
 -s_{12}c_{23}-c_{12}s_{23}s_{13}e^{\ii\delta}&
 c_{12}c_{23}-s_{12}s_{23}s_{13}e^{\ii\delta}&s_{23}c_{13}\\
 s_{12}s_{23}-c_{12}c_{23}s_{13}e^{\ii\delta}&
 -c_{12}s_{23}-s_{12}c_{23}s_{13}e^{\ii\delta}&c_{23}c_{13}
 \end{pmatrix}.
\end{equation}
\end{exercise}
\begin{solution}
官方逐元素检查重编，因子分解为编者补充。可写$V=R_{23}U_{13}R_{12}$，其中
\begin{equation}
 R_{12}=\begin{pmatrix}c_{12}&s_{12}&0\\-s_{12}&c_{12}&0\\0&0&1\end{pmatrix},\quad
 R_{23}=\begin{pmatrix}1&0&0\\0&c_{23}&s_{23}\\0&-s_{23}&c_{23}\end{pmatrix},\quad
 U_{13}=\begin{pmatrix}c_{13}&0&s_{13}e^{-\ii\delta}\\0&1&0\\-s_{13}e^{\ii\delta}&0&c_{13}\end{pmatrix}.
\end{equation}
每个因子直接由$c^2+s^2=1$得$UU^\dagger=I$，所以乘积幺正。题指定元素也可独立核查：
$(VV^\dagger)_{11}=c_{13}^2(c_{12}^2+s_{12}^2)+s_{13}^2=1$；
\begin{align*}
 (VV^\dagger)_{12}
 &=c_{12}c_{13}(-s_{12}c_{23}-c_{12}s_{23}s_{13}e^{-\ii\delta})\\
 &\quad+s_{12}c_{13}(c_{12}c_{23}-s_{12}s_{23}s_{13}e^{-\ii\delta})
   +s_{13}e^{-\ii\delta}s_{23}c_{13}=0.
\end{align*}
实项两两抵消，复项系数为$1-c_{12}^2-s_{12}^2$。分解法证明其余元素，不靠省略号假定它们成立。
\end{solution}
\begin{exercise}\label{ex:rec-14-2}
Recitation 14题2。只要CKM矩阵幺正，证明GIM机制对三代或任意代数仍成立。原题表述为“eliminating $K^0\to\mu^+\mu^-$”；顶点$u\to d+W^+$带$V_{ud}$，$d\to u+W^-$带$V_{ud}^*$。说明究竟抵消哪一项。
\end{exercise}
\begin{solution}
官方味求和证明重编并更正量词。$d,s$为不同列，$n$代幺正矩阵给
\begin{equation}
 \sum_{i=1}^{n}V_{id}^*V_{is}=(V^\dagger V)_{ds}=0.
\end{equation}
若环函数的质量无关部分为同一$F_0$，其贡献$F_0\sum_iV_{id}^*V_{is}$确为零，任意代均成立。实际环振幅可写
\begin{equation}
 \mathcal M_{s\to d}\ \propto\ \sum_iV_{id}^*V_{is}F(x_i)
 =\sum_iV_{id}^*V_{is}[F(x_i)-F(x_0)],\qquad x_i=m_i^2/M_W^2.
\end{equation}
退化上型夸克质量时整个这种环项消失；非退化质量留下差值，长程贡献也不受此一式全部消除。因此GIM抑制味变中性流及消去普适部分，不意味着物理$K^0\to\mu^+\mu^-$总衰变率严格为零。严格区分“幺正正交的列”与“质量依赖函数加权的列”，便能指出源题过强表述的问题。
\end{solution}
\originalsection{轻子对撞机中的Higgs}
\begin{exercise}\label{ex:rec-15-1}
Recitation 15题1。课堂讨论了强子对撞机$pp\to H+X$。电子--正电子对撞机如何产生Higgs？讨论两束能量对产生机制的影响。
\end{exercise}
\begin{solution}
官方机制曲线重编为自足解释和原生图，不转载其来源不明的截面图。轻子总质心能量$\sqrt s=2E_{\rm beam}$。
Higgsstrahlung为$e^+e^-\to Z^*\to ZH$，在壳阈值$\sqrt s\ge m_Z+m_H\simeq216$GeV；每束至少约108GeV，250GeV总能时两束各125GeV。其两体相空间因子
\begin{equation}
 \lambda_{ZH}^{1/2}=\sqrt{\left[1-\frac{(m_Z+m_H)^2}{s}\right]
 \left[1-\frac{(m_Z-m_H)^2}{s}\right]}
\end{equation}
在阈值为零；远高于阈值的$s$道截面大致随$1/s$下降，不能认为升高束能总会增大此道。

更高能量下，$WW$融合$e^+e^-\to\bar\nu_e\nu_eH$的重要性上升，$ZZ$融合给$e^+e^-H$。这两类交换玻色子可以为虚粒子，阈值不应写成$2m_W+m_H$；高能有效玻色子近似显示融合截面随对数缓慢增长。再高能可研究$e^+e^-\to t\bar tH$，在壳阈值$2m_t+m_H$，约470GeV，直接探测top Yukawa。直接$e^+e^-\to H$的电子Yukawa耦合正比$m_e$，通常极小。
\begin{center}\begin{tikzpicture}[scale=.62]
\draw[fermion](-1.4,-1)--(0,0);\draw[fermion](0,0)--(-1.4,1);
\draw[photon](0,0)--(1.5,0)node[midway,above]{$Z^*$};
\draw[photon](1.5,0)--(3,1)node[right]{$Z$};\draw[dashed](1.5,0)--(3,-1)node[right]{$H$};
\node at(.8,-1.8){Higgsstrahlung};
\draw[fermion](5,1.5)--(6,1)--(8,1.5)node[right]{$\nu_e$};
\draw[fermion](8,-1.5)--(6,-1)--(5,-1.5);
\draw[photon](6,1)--(7,0)node[midway,above]{$W$};\draw[photon](6,-1)--(7,0);
\draw[dashed](7,0)--(9,0)node[right]{$H$};
\node at(6.5,-1.8){$WW$融合};
\end{tikzpicture}\end{center}
图右上下轻子线分别连接相应束流与中微子末态；$ZZ$融合只须将交换线及末态改为中性$Z$和散射电子对。比较这些道时，还要声明束流极化、实验接受度和辐射修正；本题未给这些数值，不能从原示意图虚构精密截面。
\end{solution}
